\section{The endpoint and the sampling law}

Let \(R(n)\) be the number of ordered pairs of primes with sum \(n\). The
Hardy--Littlewood main term for \(R(n)\), for even \(n\), is
\(n\SSer(n)/(\log n)^2\), where
\begin{equation}\label{eq:series}
 C_2=\prod_{p>2}\left(1-\frac{1}{(p-1)^2}\right),\qquad
 \SSer(n)=2C_2\prod_{\substack{p\mid n\\p>2}}\frac{p-1}{p-2}.
\end{equation}
Along a progression with step divisible by all small primes,
the corresponding factors in \eqref{eq:series} stay fixed. The larger
prime divisors depend on the position in the progression. We determine
how large their combined effect becomes among a growing number of actual
Goldbach counts.

All logarithms are natural. Let \(N\) tend to infinity through the even
integers. Write
\begin{equation}\label{eq:parameters}
 L=\log N,\quad w=\log L,\quad W=\prod_{p\le w}p,\quad
 K=\lfloor L^D\rfloor,\quad
 X_N=\frac{N}{L^2}\,2C_2\prod_{\substack{2<p\le w\\p\mid N}}
                                 \frac{p-1}{p-2},
\end{equation}
where \(D>0\) is fixed. Set \(\sigma=8/33\). Fix
\(0<\varepsilon<(1-\sigma)/2\) and choose
\(N^{\sigma+\varepsilon}\le H\le N^{1-\varepsilon}\). Put \(Y=H/W\).
Choose \(k\) uniformly from the integers in \((Y,2Y]\), or uniformly
from the primes in that interval. In both cases let \(h=Wk\). The
probability and expectation below refer solely to this choice of \(k\).
In particular, all estimates are uniform over the even endpoints \(N\)
and admissible \(H\). For \(1\le i\le K\), put
\[
 n_i=N-ih,\qquad Z_i=(\log w)\log\frac{R(n_i)}{X_N}.
\]
Set \(Z_i=-\infty\) if \(R(n_i)=0\).

\begin{theorem}\label{thm:main}
For either sampling law,
\[
  \min_{1\le i\le K}Z_i \ \xrightarrow{\ \mathrm{P}\ }\ 0,
  \qquad
  \max_{1\le i\le K}Z_i \ \xrightarrow{\ \mathrm{P}\ }\ \log(1+D).
\]
For \(0<t<\log(1+D)\), let \(C_t=\#\{1\le i\le K:Z_i\ge t\}\).
Then
\begin{equation}\label{eq:profile}
        C_t=L^{D-(e^t-1)+o_{\mathrm P}(1)} .
\end{equation}
Precisely, for every fixed \(\eta>0\), the probability that
\(L^{D+1-e^t-\eta}\le C_t\le L^{D+1-e^t+\eta}\) tends to one.
\end{theorem}

For example, take \(K=\lfloor\log N\rfloor\). The largest value of
\(Z_i\) then tends in probability to \(\log 2\), and
\(L^{1/2+o_{\mathrm P}(1)}\) positions reach the level \(\log(3/2)\).
The endpoint \(N\) may be chosen before the step. The theorem concerns the
positive indices \(i\ge1\) and leaves \(R(N)\) undetermined.

Erd\H{o}s and K\'atai studied maxima of a nonnegative strongly additive
function along consecutive integers with the starting integer varying.
Their Theorem~4 treats the prime values \(g(p)=1/p\)
\cite{EK}. Weingartner estimated tails of the limiting distributions
of \(\sigma(n)/n\) and \(n/\varphi(n)\) by analysing high moments of
Euler products \cite{Wei}. Languasco computed fixed-order moments of
the full binary Goldbach singular series over \(n\le x\) and applied them
to Goldbach numbers in short intervals \cite{Lan}.
Here the cutoff \(w\) varies with the endpoint,
the step supplies the randomness, and all positions share this step. The
proof retains the factors at primes dividing both the index \(i\) and
the prescribed endpoint \(N\). These factors can be large for an
individual \(i\). Averaging them over the index range preserves the
exponential rate.

\section{Transferring the averaged Goldbach theorem}

Let \(\Lambda\) be the von Mangoldt function and write
\(G(n)=\sum_{p+q=n}\log p\log q\), the sum running over ordered prime
pairs. We use Theorem~1.3(iv) of Matom\"aki, Radziwi{\l}{\l} and Tao
\cite{MRT}. For fixed \(\varepsilon'>0\) in the range of that theorem,
\(x^{\sigma+\varepsilon'}\le Z\le x^{1-\varepsilon'}\), and any
\(E>0\), one has
\begin{equation}\label{eq:mrt}
  \sum_{a+b=n}\Lambda(a)\Lambda(b)
     = n\SSer(n)+O_{E,\varepsilon'}(x/\log^E x)
\end{equation}
for all but \(O_{E,\varepsilon'}(Z/\log^E x)\) even integers in
\([x,x+Z]\). This is the analytic input about binary Goldbach counts.

\begin{lemma}\label{lem:transfer}
Let \(y=L^2\), and put
\[
        T_i(k)=\sum_{\substack{w<p\le y\\p\mid n_i}}\frac1p .
\]
For each fixed \(A>0\), outside a set of \(k\) of probability
\(O_{A,D,\varepsilon}(L^{-A})\) under either law, all \(R(n_i)>0\) and
\begin{equation}\label{eq:transfer}
        \max_{i\le K}|Z_i-(\log w)T_i(k)|=o(1).
\end{equation}
\end{lemma}

\begin{proof}
For a fixed \(i\le K\), apply \eqref{eq:mrt} with
\(x=N-2iH\), \(Z=iH\) and \(\varepsilon'=\varepsilon/2\).
Here \(x\sim N\), and the factor \(i\le L^D\) fits inside the slack
at the upper end of the permissible range. Each exceptional \(n_i\)
determines at most one \(k\). Summing over \(i\) costs
\(O_{E,D,\varepsilon}(K^2H/L^E)\) choices. The prime number theorem
gives \(\log W=(1+o(1))w\). The two sampling sets have sizes
\(\asymp Y\) and \(\asymp Y/\log Y\), respectively. Choosing \(E\)
large compared with \(A\) and \(D\) proves the probability estimate.

There are \(O(\sqrt N L)\) nontrivial prime powers up to \(N\).
Their contribution for a given \(n_i\) is \(O(\sqrt N L^3)\).
For \(p,q>N/L^4\), the weight \((\log p)(\log q)\) equals
\(L^2(1+O(w/L))\). At most \(2N/L^4+O(1)\) ordered pairs have a smaller
endpoint. Since \(\SSer(n_i)\ge2C_2\), on the retained steps these estimates
give, uniformly in \(i\),
\[
 R(n_i)=\frac{N\SSer(n_i)}{L^2}
             \left(1+O\left(\frac{w}{L}+\frac{KH}{N}\right)\right).
\]
For \(p\le w\), membership in the prime divisors of \(n_i\) is fixed by
\(N\). Above \(w\), replacing
\(\log((p-1)/(p-2))\) by \(1/p\) costs \(O(\sum_{m>w}m^{-2})=O(1/w)\).
Finally \(\sum_{p>y,\,p\mid n_i}1/p\le L/(y\log y)\). Multiplication
by \(\log w\) sends all these errors to zero.
\end{proof}
